% Generated draft: requires mathematical and source review.
\documentclass[11pt]{article}
\usepackage[margin=25mm]{geometry}
\usepackage{fontspec}
\setmainfont{Latin Modern Roman}
\usepackage{amsmath,amssymb,mathtools,unicode-math}
\setmathfont{Latin Modern Math}
\usepackage{enumitem,needspace,xurl,hyperref}
\setlength{\parindent}{0pt}
\setlength{\parskip}{5pt}
\setlength{\emergencystretch}{3em}
\newcommand{\N}{\mathbb{N}}
\newcommand{\Z}{\mathbb{Z}}
\newcommand{\Q}{\mathbb{Q}}
\newcommand{\R}{\mathbb{R}}
\newcommand{\C}{\mathbb{C}}
\newcommand{\F}{\mathbb{F}}
\newcommand{\A}{\mathbb{A}}
\newcommand{\PP}{\mathbb{P}}
\newcommand{\E}{\mathbb{E}}
\newcommand{\eps}{\varepsilon}
\newcommand{\dd}{\,\mathrm{d}}
\begin{document}
\section*{UnsolvedMath ID: 700481}
\textbf{Source problem code:} catalogue entry
\subsection{Definitions and mathematical statement}
For a group \(G\), its centre is
\[
Z(G)=\{g\in G: gx=xg\text{ for every }x\in G\}.
\]
The group \(G\) is locally nilpotent if every finitely generated subgroup of \(G\) is nilpotent. Here a group \(H\) is nilpotent if its lower central series, defined by \(\gamma_1(H)=H\) and \(\gamma_{r+1}(H)=[\gamma_r(H),H]\), reaches the identity subgroup after finitely many steps; \([A,B]\) denotes the subgroup generated by the commutators \(a^{-1}b^{-1}ab\) with \(a\in A\) and \(b\in B\).

The supplied question asks whether there exists a locally nilpotent group \(G\) such that \(Z(G)=\{1\}\) and \(G\) satisfies the \emph{weak maximal condition for normal subgroups}.

The record does not define this last condition. A candidate interpretation, requiring confirmation against the source, is that there is no infinite ascending sequence
\[
N_0<N_1<N_2<\cdots,\qquad N_i\trianglelefteq G,
\]
with \([N_{i+1}:N_i]\) infinite for every \(i\geq 0\). Under this interpretation, the question is
\[
\exists\,G\quad
\begin{cases}
\text{every finitely generated subgroup of }G\text{ is nilpotent},\\
Z(G)=\{1\},\\
\text{no infinite ascending sequence of normal subgroups of }G\\
\text{has all successive indices infinite}.
\end{cases}
\]
This is a conditional formalization, pending verification of the named chain condition. The source also does not explicitly require \(G\neq\{1\}\); whether nontriviality is intended must be checked before adopting a definitive statement.
\subsection{Short English statement}
L. A. Kurdachenko asks whether a locally nilpotent group can have trivial centre while satisfying the weak maximal condition for normal subgroups. The supplied record leaves the chain condition undefined and does not explicitly require the group to be nontrivial.
\subsection{Sources}
\begin{itemize}
\item Supplied local UnsolvedMath registry record: permanent ID \(700481\), exact source code \texttt{catalogue entry}, title \emph{Kourovka Notebook Problem 8.29}; upstream revision \texttt{d371ee917ac9647fbb06cd6fc529deec6f3dda82}.
\item Attributed in the supplied record to L. A. Kurdachenko, proposed in 1982; original problem number 8.29 in \emph{The Kourovka Notebook (Unsolved Problems in Group Theory)}, 8th issue (1982).
\item The supplied background also identifies the 21st issue edition (2026), arXiv:\texttt{1401.0300v48}. The collection description names E. I. Khukhro and V. D. Mazurov as editors. These bibliographic details are supplied attribution and have not been independently verified here.
\item The registry labels the problem \texttt{open}. This is a supplied status label, not an independently verified literature assessment. No source excerpt defining the weak maximal condition is supplied.
\end{itemize}
\end{document}
